\documentclass[pdflatex,sn-mathphys-ay]{sn-jnl}

\usepackage{graphicx}
\usepackage{amsmath,amssymb}
\usepackage{booktabs}
\usepackage{xcolor}

% Marks text still to be written; the final paper must contain none.
\newcommand{\placeholder}[1]{\textcolor{red}{[#1]}}

\raggedbottom

\begin{document}

\title[In RAG We Trust?]{In RAG We Trust? A Controlled Measurement of How Open
Language Models Handle Poisoned Retrieval}

\author*[1]{\fnm{Angela} \sur{Aldaoud}}\email{angela.aldaoud@studenti.unimi.it}

\affil*[1]{\orgdiv{M.Sc.\ in Computer Science (LM-18)},
\orgname{Universit\`a degli Studi di Milano}, \orgaddress{\city{Milan},
\country{Italy}}}

\abstract{\placeholder{Abstract: written last, after the results.}}

\keywords{retrieval-augmented generation, knowledge conflicts, data poisoning,
robustness, evaluation of language models}

\maketitle
% ---------------- PAGE BUDGET (max 8 pages incl. floats and references) ----
% Title, abstract, keywords ............ 0.50
% 1 Introduction (incl. literature) .... 1.00
% 2 Research Question and Methodology .. 2.50  (incl. revisions table)
% 3 Experimental Results ............... 2.50  (setup 0.25; RQ1-RQ4 2.25)
% 4 Concluding Remarks ................. 0.75
% Back matter (AI use, availability) ... 0.25
% References ........................... 0.50
% Total ................................ 8.00
% Floats: max 4 (revisions table + 3 in Results).
% Overflow: cut Methodology detail (refer to repo docs), never the discussion.
% ---------------------------------------------------------------------------


\section{Introduction}\label{sec:intro}
\placeholder{Motivation and the RAG trust problem.}

\placeholder{Short discussion of the relevant literature.}

\placeholder{Main question and contribution.}

\section{Research Question and Methodology}\label{sec:method}

\subsection{Goals and research questions}\label{sec:rqs}
\placeholder{Main question, RQ1--RQ4, hypotheses H1--H4.}

\subsection{Problem definition}\label{sec:formal}
\placeholder{Formal notation, outcomes, metrics.}

\subsection{Data and poisoned contexts}\label{sec:data}
\placeholder{HotpotQA sample, false answers, poison types, context rule,
retrieval check.}

\subsection{Models, prompts and decoding}\label{sec:models}
\placeholder{Three models, three prompts, greedy decoding and samples.}

\subsection{Scoring and statistical analysis}\label{sec:scoring}
\placeholder{Outcome rules, conflict-flag judge and its validation,
hallucination, statistics.}

\subsection{Revisions to the planned design}\label{sec:revisions}
\placeholder{What changed after testing, and why.}

\section{Experimental Results}\label{sec:results}

\subsection{Experimental setup}\label{sec:setup}
\placeholder{Written in Block G2.}

\subsection{Robustness to poison (RQ1)}\label{sec:rq1}
\placeholder{Written in Block G2.}

\subsection{Effect of meta-prompts (RQ2)}\label{sec:rq2}
\placeholder{Written in Block G2.}

\subsection{Conflict awareness (RQ3)}\label{sec:rq3}
\placeholder{Written in Block G2.}

\subsection{Poison types, models and realistic retrieval (RQ4)}\label{sec:rq4}
\placeholder{Written in Block G2.}

\section{Concluding Remarks}\label{sec:conclusions}
\placeholder{Critical discussion, limitations and future work, written in
Block G2.}

\backmatter

\bmhead{Use of AI tools}
\placeholder{Mandatory AI-usage disclaimer, written in Block G2 from the
AI-usage log.}

\bmhead{Code and data availability}
\placeholder{GitHub repository link and how to obtain the raw outputs
(Block H).}

\bibliography{refs}

\end{document}